\section{Initiatic Centers and History (I)}

\begin{quotex}
This essay is chapter 17 of \textit{L'Arco and la Clava} by \textbf{Julius Evola}. It will be published in two parts; the first part deals with general considerations and the second with Tibet.

There are several things to ponder. For example, what is the ``symbolic dimension" to consider, and how is it related to the historical? Is there really any historical evidence, in the current academic sense, for these centers? How can such centers be recognized? What are the higher states of consciousness? Why were the centers unable to prevent the current decline, a topic Evola brings up, blaming the fall on man's freedom? Are the current pan-Islamic movements really an attempt to restore Tradition to those lands? If so, have the Sufi centers played any role in those movements?

\end{quotex}
Given the confusions that abound in this area, it is opportune to clarify first what we mean, in general, by ``initiatic centers" and ``initiatic organizations".

We treated initiation in an earlier chapter\footnote{\url{http://thompkins_cariou.tripod.com/id49.html}}, so here we will limit ourselves to recall that, in its authentic and integral meaning, initiation consists in an opening of consciousness beyond human and individual conditionalities, entailing a modification of the subject (of his ``ontological status"), who then participates in a higher freedom and a higher consciousness. That is tied in the grafting in the individual of an influence in a certain transcendental, or not simply human, way. In general such an influence is transmitted, and the transmission is an essential function of an initiatic center. What arises from it is the idea of an uninterrupted ``chain" (the term used in Islam is exactly silsila) whose origins are remote and mysterious, in parallel with a ``tradition". According to the Guenonian school, the various initiatic centers, if they are authentic and regular, would be connected to a unique center and would even draw its origins from it. A similar assumption confronts us however with difficult problems, if it must be validated in some way.

For the topic that we intend to discuss here, the question arises about that aspect of spiritual influences that do not solely concern ``knowledge", a spiritual illumination, the attainment of gnosis, but also a power. This power could even be considered by some—and not incorrectly—as a positive sign because, as long as it is only about a knowledge regarding higher spheres, but if it remains in a purely interior domain, it could also create illusions. The coexistence of a power, that as such is verifiable, is an indirect, but sufficiently positive, proof, for the solidity and the reality of the same knowledge which is considered to having been attained by means of an initiation.

Titus Burckhardt, moreover, was able to speak, in terms of initiation centers, of spiritual influences ``whose action, if not always apparent, passes incommensurably beyond everything that is in the power of men". Let's move now into the area of reality and history. We had a friendly debate with Burckhardt about the existence and the state of initiatic organization in the world of today. It is not that we assert that they no longer exist, but rather that they have become ever more rare and difficult to access (the assumption is always that it is about authentic initiatic organizations, and not of spurious groups that claim to have such a character). It seems that a progressive withdrawal of such organizations, and therefore of the forces that are manifested in them and of which they were the bearers, has been verified. Moreover, to refer to some traditions worthy of note, this phenomenon would not even be recent. We will limit ourselves to mention those texts in which it is said that the search for the Graal was brought to an end, but through a divine order, the Templars of the Graal abandoned the West and moved, together with the mystical and magical object that no longer remained ``among the sinful peoples", to a mysterious region, sometimes identified as where ``Prester John" reigned. And there, magically, the castle of the Graal, Montsalvat would also have been relocated. Naturally, the symbolic dimension must be considered in all this.

A second, more recent tradition concerns the Rosicrucians. After causing a commotion, especially with their Manifestos in which they made known their ``invisible and visible presence" and with their project of the restoration of a general higher order, the Rosicrucians also withdrew, going back to the beginning of the 18th century; for this reason, we consider that certain groups who subsequently self-qualified themselves as ``Rosicrucian" were without authorization and lacking every regular traditional filiation or continuity.

We could add an Islamic testimony of the Ismailist initiatic stream and in particular by that of the so-called Twelvers. The corresponding view is that the Imam, the supreme head of the Order, the manifestation of a power from above the also the beginning of initiations, has likewise withdrawn. They instead await those who will manifest again, but the current epoch would be that of an ``absence".

Nevertheless, that, in our opinion, does not imply that initiatic centers in the strict sense do not currently exist. Without doubt, they still exist, even if in that respect the West barely comes in question and needs to advert to other areas, both in the Islamic world and in the Orient. That established, the problem is the following: if, as Burckhardt asserts, such centers were deposits of the spiritual influences by definition, were deposited, apart from the initiatory use, we have to attribute the beginning of a possible action to something external that ``if not always apparent, passes incommensurably beyond everything that is in the power of men", how must we conceive the relationship between such still existent centers (if existing not as mere survivals) and the course of recent history?

From the traditional point of view, this course has, in general, an absolutely involutive and dissolutive character. Now, in the face of forces that are at work in these developments, what is the position of initiatic centers? If they always have had those influences as is claimed, must one think of a type of order received from them, but not using them, and not preventing the process of involution, or must one hold that the general process of ``solidification" and impermeability of the environment to the supersensible, provoking a type of rift now renders relative every action that goes beyond the initiatic field in the purely spiritual and interior direction?

It is good to clarify and to put aside the cases in which historically are only the harvested fruits which had been sown earlier. Men have a fundamental freedom. If they have used it for their ruin, the responsibility falls on them and there is no reason to intervene. Now, we can say that for the West, which has been taking the path of the anti-tradition for some time and that through a chain of cause and effect, sometimes quite visible, sometimes hidden to a superficial glance, fatally proceeded to find itself in its current state, that resembles the kali yuga, the ``dark age" prophesied by ancient traditions.

But in other cases, things do not remain in the same way. There are civilizations that by not having followed the same path, by not having chosen mistaken vocations, but finding themselves subjected to external influences, would have to be defended. But that seems to not be observed. For example, in the case of Islam there are certainly existing initiatic Sufi centers, but their presence has not actually prevented the evolution of Arabic countries in an anti-traditional, progressive, and modernist direction, with all the inevitable consequences.

\begin{quotex}
This is the conclusion of the essay, dealing with initiatic centers in Tibet. In it, \textbf{Julius Evola} proposed a sort of occult drone warfare, assuming initiatic centers truly possess superhuman powers. He offers reasons why the Tibetans allowed the communist invasion had they been able to prevent it by paranormal means. I am curious about reactions to this piece as well as how it relates to the current obsession with various ``superheros". 

\end{quotex}
But a decisive case is that of Tibet. Tibet did not consider at all taking the same path as Western countries. It had maintained its traditional structures intact and was also considered as a country in which, more than any other, individuals and groups existed who were in contact with super-extensible and divine powers. That did not prevent it from being invaded, profaned and devastated by the Chinese communist hordes, which also put an end to the ``myth" of Tibet whose fascination had such a hold on the Western spiritual milieu. Yet, in principle, there should have been manifested and presumed for possible use, some concrete opportunities of what was attributed to influences of a not simply human and material order.

To be precise, we are not thinking of invisible and magical barriers of protection that would have blocked the invaders of Tibet. It is sufficient to advert to something somewhat less spectacular. For example, with reference to so-called modern parapsychological research, performed under strict controls, the reality of ``paranormal phenomena" has been verified, that is, the possibility that objects can be displaced, moved or levitated at a distance without a normal explanation. Only that, given the matter with which almost exclusively has to do paranormal research, it is a question of spontaneous sporadic processes, often extrasensory, but not producible at will. Nevertheless the fact has been verified that a psychic agent can cause phenomena like the levitation of a heavy object, that implies an undoubtedly superior force to the force necessity to cause, for example, brain damage with a deadly outcome. Even the phenomena of bilocation, or of the projection of one's own image in a distant place, has been verified (moreover, it seems that it has happened even with Padre Pio of Pietralcina).

So, from everything that has been reported by travelers and observers worthy of credence, beginning with Alexandra David-Neel, similar phenomena was verified in Tibet, however not as phenomena of an extrasensory and unconscious character, but rather as consciously controlled and willed phenomena, made possible by discipline and initiations.

Now, it would have sufficed to use powers of that type to cause, for example, a cerebral lesion and thereby to strike down Mao Tse-Tung at the moment the first communist division crossed through the Tibetan border. Or rather to use that power of projection of one's own image to cause a warning apparition in the face of the Chinese communist leader.

All this should not appear as a merely digressive fantasy to those who have a conception of initiatic centers like that indicated by Burckhardt's words cited above and who believe that similar initiation centers still exist. And haven't the Tibetan traditions even spoken of the famous Milarepa who in the first period of his life, before reaching the Great Liberation, was an outlaw who dedicated himself to black magic, and, in fact, caused a massacre of his adversaries by magical means? Instead we are present at the end of Tibet, without being able to bring in the idea of a type of Nemesis (as for the West). A recently translated book in Italian\footnote{\textit{Born in Tibet}, by \textbf{Chogyam Trungpa}.} speaks of the Odyssey of those lamas who were able to do nothing but escape in order to save their lives, while in the country others were massacred, by those who sought to eradicate everything that had a holy character, and began the communist atheist indoctrination of the people. The only resistance was the guerrillas of Tibetan partisans who pulled back in an inaccessible area. It is useless to say what, an occult defense like that mentioned would have signified instead. Its significance would have made all the marvels of voyages and amazing explorations boasted about by the modern Western world appear very banal and insipid.

So the problem we just posed holds good, without, it would seem, an adequate clarification of it being possible. The only idea that perhaps I could put forward is that of a type of fracture of what exists and the autonomization of a certain part of reality, therefore also of history, with a consequent impermeability in regard to extrasensory influences. It could also relate to the doctrine of cycles, to that which is characteristic of the closing of a cycle. Only that in the case in question, little space would remain for values of a moral character. We should think of a general process in which even those who have not fomented it find themselves implicated. And it should also relate to a type of password transmitted to initiatic centers, with the goal of letting destinies be fulfilled.

This is an order of ideas that would lead rather far, to the same conception of an inscrutable direction of the world and, in another direction, to the existing relation between freedom and necessity. When another perspective should not be of value, necessity could be related only to the factual domain of existence and freedom for various attitudes that could be assumed in the face of facts (or to the reaction to them), which, in terms of principle, is not determined. In this context one could give, among others, a particular weight even to negative and dramatic experiences, that they can give if one assumes a given attitude, even to take on the character of a test. As we see, this is a field of rather vast and complex problems, with which is also grappled with in the theology of history. We have mentioned this only as a general background to which the deepening of the specific subject of this essay refers.

\paragraph{Note}
Evola added the following note regarding the theology of history.

\begin{quotex}
At the traditional Catholic base, things are not so easy in cases like Spain's so-called ``Invincible Armada"; organized against the heretics and set sail after having every very solemn consecration, it was destroyed, even prior to actual combat, by a storm, the ``forces of nature". 

\end{quotex}


\flrightit{Posted on 2013-01-20 by Aeneas }

\begin{center}* * *\end{center}

\begin{footnotesize}\begin{sffamily}



\texttt{Jason-Adam on 2013-01-21 at 08:24 said: }

Some modern writers that that Evola was maybe a bit pessimistic regarding the influence of iniatory centres in the Islamic world and that the modern Islamic movements such as the apocalyptic Khomeini movement in Iran, based around the coming return of the 12th Imam, is such a result of a traditional force acting.


\hfill

\texttt{Cologero on 2013-01-27 at 20:45 said: }

It seems Evola is moving the problem of theodicy from God to the initiatic centers. If the latter centers had been able to prevent a great evil, should they not have done so? he asks. He makes some conjectures about why that may not have been possible.

You can imagine how absurd history would be if every ``bad guy" was occultly punished by initiatic centers prior to their nefarious deeds. Given such powers on the material plane, it is difficult to conceive how they could have disappeared. Isn't it more likely that the superhuman powers spoken of involve states that transcend human, all-to-human, concerns?

His theology of history is also skewed. If God always ``intervened" to punish the bad and reward the good in this life, then men would become religious merely to be on the right side of things, in order for personal gain. There would certainly be no need for faith, since God's will would be visible to everyone, clearly and not as a reflection in a dirty mirror.


\hfill

\texttt{Logres on 2013-01-27 at 23:22 said: }

If the natural state of affairs is the ``Golden Age", then to induce a degeneration, the falling asleep of higher centers (collectively and personally) would be required: it would grow colder and colder, and conscious centers would drop off to sleep. Those that were still awake would perhaps withdraw in order to maintain their own equilibrium? Basically, a period of dormancy? This is in Nature. Maybe, as you say, the whole point is to reassert the absoluteness of God at the end of the cycle, in such a way that no one who remains awake can be in any doubt that God's sovereignty is supreme. He seems to be referring at the end of the essay to ``transvaluation of morals"; and I have no idea what he means by a ``password". But the Spanish Armada reference is intriguing, especially when one thinks that Lepanto had an entirely different outcome in the opposite direction (and not that far removed in time).


\hfill

\texttt{Cologero on 2013-01-27 at 23:59 said: }

Password or watchword … the translation of ``parola d'ordine'.

Yet Constantinople fell not long before the Armada and a millennium after Constantine's victory. Notre Dame lost to Alabama.

The rain falls on the just and the unjust … I don't think we can draw any conclusions from who gets wet.

But he was writing specifically about Tibet which in recent times (Evola's recent) was presumed to still have secret initiatic centers, or perhaps not so secret. They had the powers of telekinesis and bilocation. If true, why did they not use them against the communists? Of course, these superhuman powers existed in a vertical direction, not on the horizontal plane. In line with the principle ``don't feed the animals in the wild", and ``let wildfires run their course", these higher powers don't interfere directly on lower planes.


\hfill

\texttt{Mihai on 2013-01-28 at 05:15 said: }

This article surprises me. As an author of so many good esoteric books (and a person who studied such things all his life) Evola seems to speak with no reservations about reckless usage of magical powers, with an ease that only the lesser of occultists would be capable of. 

It surprises me even more greatly since, even as early as 1931, in his ``Hermetic tradition", he wrote about magical powers for material ends that ``a true adept would only make use of such a thing when it is justified by a higher, spiritual goal". (the quote is approximate, from my memory- see the closing chapters of the above mentioned book). 

If we are to approach historical things from such an angle, we could ask indefinetly all sorts of non-sense questions such as : why did the early Christian martyrs, who produced such great miracles, not use their powers to (at least) avoid detection from the Roman authorities when they came looking for them. Or why did Christ, as the Second Person of the Trinity and Principle of the whole of manifestation, did not use His powers to come down from the Cross (as His adversaries ``suggested"), or, better yet, to tear them to shreds, or to destroy the whole universe at that moment, as unworthy to receive His blessings.

It is very unclear to me why would Evola, who his whole life emphasized the overcoming of a merely human understanding of history, would make such inquires in this essay. 

As for the theodicy that Cologero brought forth in the above comment, I think that Boethius did a fine work.


\hfill

\texttt{Cologero on 2013-01-28 at 07:58 said: }

Not to forget the first half of the essay where Evola is addressing Burckhardt's claim about initiatic centers ``whose action, if not always apparent, passes incommensurably beyond everything that is in the power of men". Evola seems to be saying that if that is true, then their action in history is not visible. Recall this passage:

\begin{quotex}
From the traditional point of view, this course has, in general, an absolutely involutive and dissolutive character. Now, in the face of forces that are at work in these developments, what is the position of initiatic centers? If they always have had those influences as is claimed, must one think of a type of order received from them, but not using them, and not preventing the process of involution, or must one hold that the general process of ``solidification" and impermeability of the environment to the supersensible, provoking a type of rift now renders relative every action that goes beyond the initiatic field in the purely spiritual and interior direction? 

\end{quotex}

\hfill

\texttt{Mihai on 2013-01-28 at 08:54 said: }

Yes, it seems I have taken this text out of context. 

However, it remains true that Evola does treat this question rather simplistically, although he does seem, in the end, to at least lean towards a right conclusion. The problem is to understand, first of all: what do we mean by an invisible ``action". Surely, no such center would ever act against the course of Divine order. 

It is just a question of HOW do these `centers' manifest beneficent effects, according to the natural flow of each age. For example, in the current state of the Iron Age, to act upon the world may mean to lead a few of the individuals who still are capable of upholding superior principles back to tradition. 

I believe that Evola's perspective does tend, at times, to be limited because of his somewhat ``magical" understanding of certain esoteric phenomena (if I can call them such), a perspective he never really quite shook of.


\hfill

\texttt{Scardanelli on 2013-01-28 at 09:22 said: }

So the modern superhero phenomena could be seen as a sort of counter-initiate… He who uses his superhuman powers strictly on the horizontal plane for material gains and ends, and always to the benefit of the masses. For a traditional initiate to use his powers in this way would be a debasement of himself.


\hfill

\texttt{Scardanelli on 2013-01-28 at 12:16 said: }

Further, would this not also imply that physical superhuman powers are but the lower analogue of a higher understanding? For instance, telekenesis is simply the physical manifestation of a power which is manifested differently at a higher level? I'm not sure what the higher analogues might be that correspond to their lower physical powers, but perhaps this is why I can't fly through the air or read the thoughts of others.


\hfill

\texttt{Scardanelli on 2013-01-28 at 13:10 said: }

Thus, it would seem that it isn't a question of why these initiatic centers failed to act, as it would seem that in the ``solid" nature of the Kali Yuga especially, these purely physical powers might be easier or encouraged in some way. Perhaps any attempt to physically act against the forces of degeneracy would only encourage the process…you hint at this somewhat in your comment Cologero, and forgive me if I'm repeating the obvious here. I'm just following a line of thought…


\hfill

\texttt{Jason-Adam on 2013-01-28 at 23:59 said: }

Very well written and thoughtful piece by Evola here.

The only thing I can add to his writing is by asking if there is any merit at all to the hypothesis of a few so-called traditionalists that there were somehow connections between initiatory centres and communist parties and that the triumph of communism in the east was not in spite of but because of said centres ?


\hfill

\texttt{Avery on 2013-01-29 at 02:02 said: }

Evola had a bit of blurring between fantasy and reality in his most abstruse studies of distant places and ancient lands, even while his practical understanding of the world around him was solid. If you read The King of the World carefully, you will notice that even when Guénon felt personally convinced of a fact about magic, initiation, deep history, etc., he found its symbolism to be more important and make the pursuit of ``authenticity" less relevant. Evola, on the other hand, really was interested in Atlanteans and superpowered Tibetan monks, even though he saw through the myths of contemporary Italy. It's sort of a modern version of the traditional Eastern story of getting distracted by promises of unlimited power while following the initiatory path.


\hfill

\texttt{Pickman on 2013-01-29 at 14:05 said: }

[Cologero :Of course, these superhuman powers existed in a vertical direction, not on the horizontal plane. In line with the principle ``don't feed the animals in the wild", and ``let wildfires run their course", these higher powers don't interfere directly on lower planes.]

What are you suggesting that these so called superhuman powers were metaphors of in a vertical direction?


\hfill

\texttt{Cologero on 2013-01-29 at 16:51 said: }

Yes, Pickman, primarily metaphors or better, symbols. This is not to suggest that they are not ``real" powers at their own level.

Otherwise, we are stuck with the problem Evola brings up: why did the centers not use their powers over material manifestation to prevent the communist takeover? Instead, the lamas took off when they could, leaving the population to their own devices.

We see this, too, in our time with the ``Law of Attraction" fad. For its adherents, the purpose of higher spiritual powers is to manifest better things in life such as automobiles, money, love, etc. This would eliminate the whole notion of a ``quest", or seeking, since who were of superior spiritual development would be visibly obvious. But mythically, the quest is into the unknown, not the known.


\end{sffamily}\end{footnotesize}
